\documentclass[11pt]{article}
\usepackage[margin=2.6cm]{geometry}
\usepackage{mathpazo}
\usepackage{microtype}
\usepackage{booktabs}
\usepackage{graphicx}
\usepackage{float}
\usepackage{xcolor}
\usepackage{titlesec}
\usepackage{caption}
\usepackage{listings}
\usepackage[hidelinks]{hyperref}
\emergencystretch=3em

\definecolor{accent}{HTML}{1F4E79}
\definecolor{good}{HTML}{1A7F37}
\definecolor{bad}{HTML}{B42318}
\titleformat{\section}{\Large\bfseries\color{accent}}{\thesection}{0.8em}{}
\titleformat{\subsection}{\large\bfseries\color{accent!80!black}}{\thesubsection}{0.8em}{}
\captionsetup{font=small,labelfont={bf,color=accent}}
\newcommand{\armA}{\textcolor{bad}{arm A}}
\newcommand{\armB}{\textcolor{accent}{arm B}}

\lstdefinestyle{json}{
  basicstyle=\ttfamily\footnotesize,
  backgroundcolor=\color{gray!8},
  frame=single, rulecolor=\color{gray!40}, framerule=0.4pt,
  breaklines=true, columns=fullflexible, keepspaces=true,
  showstringspaces=false, aboveskip=8pt, belowskip=8pt}
\lstset{style=json}

\newsavebox{\eliBox}
\newenvironment{eli5}
  {\par\medskip\noindent\begin{lrbox}{\eliBox}\begin{minipage}{\dimexpr\linewidth-2\fboxsep-2\fboxrule}\small}
  {\end{minipage}\end{lrbox}\noindent\fcolorbox{accent!40}{accent!4}{\usebox{\eliBox}}\par\medskip}

\title{\color{accent}\bfseries Structured DAG-Diff Mutation for CARL Reasoning Chains\\
\large Weak-Model Stress A/B: does constrained decoding hold when the mutator is cheap?\\
\large (single-parent, \texttt{llama-3.1-8b}, iteration-matched on the first 100 mutants)}
\author{GigaEvo}
\date{2026-07-03}

\begin{document}
\maketitle

\begin{abstract}
\noindent
We replaced free-form JSON rewriting of CARL reasoning-chain genomes with a
\emph{schema-constrained positional-slot diff}: the mutation LLM emits a typed edit whose
grammar is compiled per parent, enforced by the provider's constrained decoding, validated
with pydantic, and applied deterministically. Two prior A/B runs on a strong reasoning
mutator (gemini-3.5-flash) showed the diff eliminating parse failures at a token premium.
This run stresses the opposite regime: we deliberately picked a \emph{weak, cheap,
non-reasoning} 8B mutator (\texttt{meta-llama/llama-3.1-8b-instruct}) to make the free-form
path fail, ran both arms from scratch with \texttt{num\_parents=1}, stopped the slow rewrite
arm at 100 accepted mutants, and compared both arms over their \textbf{first 100 mutants,
iteration-matched}. The contrast is now stark. \textbf{The rewrite arm persisted 43 invalid
children out of 100} (40 malformed-JSON, 3 other) and needed 139 LLM calls to accept 100
mutants (39 pre-persist generation failures on top); \textbf{the diff arm persisted 0 invalid
children out of 100}, took 101 calls (one retry), and logged a single benign
diff-schema rejection caught before persist. Because the model is non-reasoning
(reasoning tokens $=0$ on both arms), the diff's real output saving is finally unmasked:
\textbf{output 42.5K (diff) vs 120.4K (rewrite) --- the diff spends 35\% of the rewrite's
output tokens} (mean 419 vs 1{,}204 per completed call), at comparable input (478K vs 460K).
Fitness is at parity at the top ($\sim$0.50 both) with a higher diff child-mean (0.413 vs
0.362): constrained edits degrade less. The two arms diverge structurally --- the rewrite's
best child is a \emph{degenerate 1-step collapse} of the 4-step seed (0.506), while the diff
preserves coherent multi-step chains (91/100 at 8 steps) and its best over the full run is a
legible 4-step prune (0.525, above the rewrite's peak). We reconstruct that best diff program
as a series of slot-diffs from the seed (\S\ref{sec:series}): extend $4\!\to\!8$ steps, carry
the structure through re-evaluations, prune back to 4 --- every hop a human-readable, provably
well-formed edit.
\end{abstract}

\section{Motivation}
CARL chain genomes are wire-JSON documents (\texttt{ReasoningChain.to\_dict} plus platform
extras). The stock mutation path asks the LLM to rewrite the whole document as free text:
every emission re-risks the entire syntax, and a malformed child is only discovered
\emph{after} it has been persisted and (for a parse that survives extraction) has consumed a
full evaluation. The hypothesis: constraining mutation to a typed diff language --- where
structural validity is enforced by the provider's constrained decoding and a pydantic
contract --- eliminates that failure class at the cheapest possible stage.

The earlier A/B pair ran this on \texttt{gemini-3.5-flash}, a capable reasoning model that
rarely emits malformed JSON unprompted; there the rewrite arm lost only a handful of mutants
and the argument rested on principle (unrepresentable-by-construction) more than on a large
observed gap. The obvious question: \emph{does the guarantee matter when the mutator is
weak?} A cheap non-reasoning model is exactly where free-form structured emission is known to
break --- truncated objects, dropped required keys, prose leakage. If constrained decoding is
worth its schema cost, this is where it should show.

\section{Setup}
\label{sec:setup}
Two arms, run from scratch, identical except the mutation operator:
\begin{itemize}
  \item \textbf{Mutator:} \texttt{meta-llama/llama-3.1-8b-instruct} via OpenRouter --- cheap,
  older, \emph{non-reasoning}. Chosen after probing candidates: its provider
  \emph{enforces} the \texttt{json\_schema} contract (so the diff arm's grammar is honoured),
  yet it is weak enough to break free-form JSON emission ($\sim$25\% malformed in a probe).
  \item \textbf{Executor:} \texttt{Qwen3-235B-A22B-Instruct} via the LiteLLM proxy, scoring
  each child by mean ROUGE-L F1 on the summarizer eval set (higher better).
  \item \textbf{Config:} \texttt{problems/chains/summarizer}, \texttt{memory=none},
  \texttt{num\_parents=1}, 4 seeds, disk storage. \armA{} = stock full wire-JSON rewrite with
  a chain-faithful prompt; \armB{} = structured slot diff (\S\ref{sec:language}).
  \item \textbf{Stopping \& matching:} the rewrite arm is the slow one (it retries failed
  generations), so it was stopped at exactly 100 accepted mutants; the diff arm ran to 250+.
  \textbf{All comparisons below are iteration-matched on each arm's first 100 accepted
  mutants} (retries and pre-persist failures counted as the cost of reaching 100).
\end{itemize}

\noindent\textbf{Single parent.} With \texttt{num\_parents=1} the diff grammar collapses to
its single-parent form: \texttt{base\_parent} is one letter, and every \texttt{keep}-id enum
ranges over one parent's steps. Crossover (keeping a donor parent's step) is therefore
unavailable in this run --- a deliberate simplification, with the two-parent crossover schema
left as prior work. The grammar is otherwise byte-for-byte the shipped one; the single-parent
schema ($\sim$16.6\,KB) is within 1\% of the two-parent schema, so this does not materially
change the constrained-decoding load.

\section{The diff language}
\label{sec:language}
The child chain is a shelf of eight numbered slots; the LLM fills them left-to-right. This is
the shipped fixed-key object encoding (E4) --- the same grammar the earlier A/B ran.

\begin{eli5}
Imagine the child chain as a shelf of eight numbered boxes. To describe a new chain, you fill
boxes from the left. Into each box you put either a \emph{photocopy} of a step from the parent
--- naming it by its printed label, like \texttt{a2}, optionally with corrections scribbled on
top --- or a \emph{brand-new} step you write yourself. Every box also declares which earlier
boxes it reads from, and here is the trick: the form for box $k$ only has checkboxes for boxes
$1..k-1$. Pointing at yourself or at a later box is not ``against the rules'' --- there is
literally no place on the form to write it. Box~1 has no checkboxes at all; it reads the raw
task input. When you are done, the remaining boxes stay empty (no gaps allowed), and whatever
sits in the last filled box produces the answer that gets graded.
\end{eli5}
\noindent
The LLM never writes a chain; it fills in this form. The form is a JSON schema compiled fresh
for each mutation from the actual parent, and the provider's constrained decoding guarantees
the reply matches the form --- a violating token can never even be sampled. Alongside the
slots the diff carries an \emph{evidence surface} identical to the standard mutation
operator's (\texttt{archetype} from a shared vocabulary, a falsifiable \texttt{justification},
\texttt{insights\_used}/\texttt{insight\_ids\_used}, \texttt{card\_ids\_used},
\texttt{changes}), so diff children self-report change in the same shape as rewrite children
and flow through the citation-integrity check on the same path.

\subsection{What the grammar makes impossible}
The design goal: a defect class should be \emph{unrepresentable}, not detected. What an
unconstrained rewriter routinely gets wrong, this language cannot even spell:
\begin{table}[H]\centering\small
\begin{tabular}{ll}
\toprule
illegal child & why it cannot be emitted \\
\midrule
self-dependency (\texttt{slot\_k} $\to$ \texttt{slot\_k}) & slot $k$'s enum stops at \texttt{slot\_{k-1}} \\
forward dependency (\texttt{slot\_2} $\to$ \texttt{slot\_5}) & same enum narrowing \\
any dependency on slot 1 & the field does not exist there; extra keys rejected \\
dependency cycle & impossible transitively: all edges point backwards \\
keep of a nonexistent step (\texttt{a9}) & \texttt{id} enum is compiled from the actual parent \\
9-step chain & there is no \texttt{slot\_9} property \\
empty chain & \texttt{slot\_1} is required \\
empty \texttt{title}/\texttt{aim} & \texttt{minLength: 1} in the schema \\
runaway \texttt{dependencies} list & \texttt{maxItems}: $k-1$ caps the array \\
invented fields, malformed JSON & \texttt{additionalProperties: false}; constrained decoding \\
\midrule
gap fill (\texttt{slot\_3} filled, \texttt{slot\_2} null) & \emph{the} residual pydantic validator (contiguity) \\
\bottomrule
\end{tabular}
\caption{Everything above the line is enforced by the decoding grammar itself. Contiguity is
the single rule left to a validator. Enforcement is a two-layer contract: per-mutation schema
compilation (\texttt{AllowedDagChanges.build\_schema}) shipped to the provider as the
\texttt{json\_schema} structured-output grammar, then the same pydantic model as a backstop
before anything touches the engine. A rejection costs one LLM call --- nothing persisted,
nothing evaluated.}
\end{table}

\subsection{Operator and vocabulary seam}
Two pieces, both behind existing seams:
\path{gigaevo/evolution/mutation/structured_diff.py} is the generic engine-level operator
\texttt{StructuredDiffMutationOperator} (ask an \texttt{AllowedChanges} vocabulary for a
\texttt{DiffSchema}, run a structured-output call, validate, apply --- it knows nothing about
chains); \path{gigaevo/chains/dag_changes.py} is \texttt{AllowedDagChanges}, the chain-specific
vocabulary --- schema compilation, parent rendering (the \texttt{a1}/\texttt{a2} listings the
LLM sees), the deterministic applier, and the natural-language contract injected into the
prompt. The applier walks \texttt{slot\_1..slot\_8} to the first \texttt{null}, copies content
from the referenced parent step and overlays \texttt{edits}, turns \texttt{dependencies} into
sorted integer positions, stamps the last step as the output, and round-trips through
\texttt{ReasoningChain} as a final assertion --- the child is indistinguishable in format from
a hand-written genome.

\section{Results}
\label{sec:results}

\subsection{Validity funnel --- the headline}
\begin{table}[H]\centering
\begin{tabular}{lcc}
\toprule
first 100 accepted mutants & \armA{} (rewrite) & \armB{} (diff) \\
\midrule
mutation LLM calls to accept 100 & 139 & \textbf{101} \\
\quad calls per accepted mutant & 1.39$\times$ & \textbf{1.01$\times$} \\
pre-persist generation failures & 39 (no code extracted) & \textbf{1} (diff schema) \\
children persisted & 100 & 100 \\
\quad valid & 57 & \textbf{100} \\
\quad \textbf{invalid persisted} & \textbf{\textcolor{bad}{43}} & \textbf{\textcolor{good}{0}} \\
\qquad of which \texttt{json\_parse\_error} & 40 & --- \\
\qquad other validation & 3 & --- \\
\textbf{valid yield among persisted} & \textbf{\textcolor{bad}{57\%}} & \textbf{\textcolor{good}{100\%}} \\
\bottomrule
\end{tabular}
\caption{The weak mutator breaks the free-form path exactly where the design predicts. \armA{}
fails at two stages: 39 emissions never yield extractable code (retried), and \textbf{43 of
the 100 that do persist are structurally invalid} --- 40 malformed JSON (\texttt{json\_parse\_error})
that each burned an evaluation before being marked invalid, plus 3 other validation failures.
\armB{} took \textbf{101 calls for 100 mutants (one retry)} and persisted
\textbf{zero} invalid children; its one logged rejection was a diff-schema check that
fired \emph{before} persist, so nothing invalid ever entered the population. Everything the
grammar forbids (\S\ref{sec:language}) simply never arrived.}
\end{table}
\begin{figure}[H]\centering
\includegraphics[width=\textwidth]{ab_overview.png}
\caption{Left: per-child fitness (dots) and best-so-far (line); \armA's invalid children are
drawn as crosses at $y=0$ --- 40 of them, versus none for \armB. Right: the validity funnel,
first 100 accepted mutants. \armA{} needs 139 attempts and keeps 57 valid; \armB{} needs 101
and keeps 100.}
\end{figure}

\subsection{Token cost --- the diff's output saving, finally unmasked}
The earlier strong-model A/B could not cleanly show the diff's output economy because
\texttt{gemini-3.5-flash} burned a large, variable \emph{reasoning} channel that dominated the
output bill. This mutator is non-reasoning: \textbf{reasoning tokens are 0 on both arms}, so
billed output is the actual mutation payload and nothing else.
\begin{table}[H]\centering
\begin{tabular}{lccc}
\toprule
mutation side (\texttt{llama-3.1-8b}), first 100 accepted & \armA{} & \armB{} & \armB{} vs \armA{} \\
\midrule
tokens in & 460{,}446 & 483{,}088 & $+4.9\%$ \\
\textbf{tokens out} & \textbf{120{,}441} & \textbf{42{,}463} & \textbf{\textcolor{good}{$-65\%$}} \\
tokens reasoning & 0 & 0 & --- \\
mean output per completed call & 1{,}204 & \textbf{420} & \textbf{\textcolor{good}{$-65\%$}} \\
LLM calls (incl.\ retries) & 139 & 101 & $-27\%$ \\
executor completion tokens (children) & 56{,}609 & 265{,}602 & $+369\%$ \\
\bottomrule
\end{tabular}
\caption{On a non-reasoning mutator the diff spends \textbf{35\% of the rewrite's output
tokens} --- 420 vs 1{,}204 per completed call --- because it emits slot deltas, not a whole
re-serialised genome. Input is near-parity ($+4.9\%$): the diff schema is larger, but
\texttt{num\_parents=1} renders only one parent, and the schema premium that dominated the
two-parent run is offset here. \armA{} also pays 39 extra calls in retries. The one place the
diff costs \emph{more} is downstream: it preserves long chains, so its children run longer on
the executor ($4.7\times$ the completion tokens) --- the flip side of \S\ref{sec:struct}. Net
on the mutation side: the diff is both more reliable and cheaper to emit.}
\end{table}
\begin{figure}[H]\centering
\includegraphics[width=\textwidth]{ab_tokens.png}
\caption{Per-mutant token spend. Left: output tokens per mutation call (dots) with a 10-call
rolling mean; \armA{} (red) scatters higher and runs to 139 calls, \armB{} (blue) sits low and
tight and stops at 100. Right: cumulative tokens --- the output curves separate sharply
(\armB{} $\approx$ 42K vs \armA{} $\approx$ 120K) while the input curves nearly coincide.}
\end{figure}

\subsection{Fitness --- parity at the top, higher diff mean}
\begin{table}[H]\centering
\begin{tabular}{lcc}
\toprule
first 100 accepted mutants & \armA{} (rewrite) & \armB{} (diff) \\
\midrule
best child & \textbf{0.506} & 0.487 \\
child mean & 0.362 & \textbf{0.413} \\
child median & 0.385 & \textbf{0.429} \\
seed best (executor re-score) & 0.482 & 0.505 \\
best child over the \emph{full} run & (stopped at 100) & \textbf{0.525} @ iter 190 \\
\bottomrule
\end{tabular}
\caption{Both arms top out near 0.50. The rewrite owns the first-100 peak (0.506) --- but that
child is a degenerate 1-step collapse (\S\ref{sec:struct}). The diff's child \emph{mean} is
higher (0.413 vs 0.362): small schema-constrained edits stay closer to a working parent, so
the child distribution degrades less. The 4 shared seeds re-score differently across the two
runs (0.482 vs 0.505) purely from executor sampling noise; there is no same-config variance
floor for this task yet, so with one replicate per arm we claim \emph{parity}, not a fitness
win for either. The diff's best over its full run (0.525) does clear the rewrite's peak, on a
coherent 4-step chain rather than a collapsed one.}
\end{table}

\subsection{Structural exploration --- the arms diverge}
\label{sec:struct}
This is where the operators behave differently. The free-form rewriter, asked to re-emit the
whole genome, tends to \emph{simplify}: its valid children concentrate at 1--3 steps
($\{1\!:\!3,\ 2\!:\!21,\ 3\!:\!25,\ 4\!:\!6,\ 5\!:\!2\}$), and its best child is a
\textbf{single step} --- the 4-step seed pipeline collapsed to one ``read-and-summarise''
instruction. The diff operator, editing locally around an 8-step parent, keeps structure:
\textbf{91 of 100 valid children have 8 steps}. The language admits deletion (just omit a
slot), and the diff \emph{does} prune when it pays --- its own best program prunes $8\!\to\!4$
(\S\ref{sec:series}) --- but it does not collapse the chain wholesale the way an
unconstrained rewrite does.
\begin{figure}[H]\centering
\includegraphics[width=0.72\textwidth]{ab_nsteps.png}
\caption{Chain-length distribution of valid children. \armA{} (rewrite) collapses toward short
chains; \armB{} (diff) preserves the parent's 8-step structure.}
\end{figure}

\subsection{Evidence surface}
Every \armB{} child carries the aligned self-report: \texttt{justification} on 98/100,
\texttt{insights\_used} on 96/100 (parent structural insights --- \texttt{card\_ids\_used} is
0 everywhere, consistent with \texttt{memory=none}), a \texttt{changes} list on 32/100, and
\texttt{archetype} drawn from the shared vocabulary --- though the weak model uses only two of
the six labels (Guided Innovation 52, Precision Optimization 48), against the six-way spread
the strong model produced. The instrumentation is live; a memory-on run would give the
citation-integrity check something to score.

\section{Best programs, and the diff series that built one}
\label{sec:series}
\subsection{The two arms' best children look nothing alike}
\armA's best child (iter 59, fitness 0.506) is a \textbf{one-step chain} --- the rewriter
threw away the seed's Read $\to$ Select $\to$ Draft $\to$ Finalise pipeline and emitted a
single instruction:
\begin{lstlisting}
step 1 (output): "Summarize paragraph"
  aim: Produce exactly one concise sentence that captures the main idea.
  action: Read the source text from the Data section; reply with one sentence.
\end{lstlisting}
It scores well on ROUGE-L, but the DAG is gone: there is nothing left to mutate structurally,
and no intermediate reasoning to inspect. This is the free-form failure mode that is not a
crash --- a \emph{valid} child that has quietly discarded the search space.

\subsection{\armB's best, reconstructed as slot-diffs from the seed}
\armB's best program over the full run (iter 190, fitness 0.525) keeps a coherent 4-step
chain, and --- because every mutation is a typed diff --- its entire history is legible. The
single-parent lineage from the seed is six nodes; each hop is a well-formed, human-readable
edit (the grammar guaranteed it before it was ever evaluated):
\begin{center}\small
\begin{tabular}{llll}
\toprule
node & fitness & steps & the diff (archetype; move) \\
\midrule
seed & 0.432 & 4 & --- (Read $\to$ Select $\to$ Draft $\to$ Finalise) \\
gen2 & 0.436 & 8 & Guided Innovation: \textbf{extend $4\!\to\!8$}, edit 4 stage-actions/aims \\
gen3 & 0.487 & 8 & Precision Optimization: keep all 8 (re-eval) \\
gen4 & 0.421 & 8 & Precision Optimization: keep all 8 (re-eval) \\
gen5 & 0.506 & 8 & Precision Optimization: keep all 8 (re-eval) \\
\textbf{gen6} & \textbf{0.525} & \textbf{4} & Precision Optimization: \textbf{prune $8\!\to\!4$}, keep slots 1--4 \\
\bottomrule
\end{tabular}
\end{center}
The two \emph{structural} hops are the interesting ones, and both are one-line diffs. The
extension (gen2) duplicates parent steps into new slots and rewrites their content:
\begin{lstlisting}
{"archetype": "Guided Innovation", "base_parent": "A",
 "slot_1": {"id": "a1", "kind": "keep",
            "edits": {"stage_action": "Use a cached or pre-loaded source text"}},
 "slot_2": {"id": "a2", "kind": "keep",
            "edits": {"aim": "Segment the text into claims before selecting the main idea"}},
 "slot_3": {"id": "a3", "kind": "keep"},  "slot_4": {"id": "a4", "kind": "keep"},
 "slot_5": {"id": "a3", "kind": "keep",
            "edits": {"stage_action": "Polish the drafted sentence into one sentence"}},
 "slot_6": {"id": "a4", "kind": "keep"},
 "slot_7": {"id": "a4", "kind": "keep",
            "edits": {"stage_action": "Maintain language consistency"}},
 "slot_8": {"id": "a4", "kind": "keep"}}
\end{lstlisting}
The prune (gen6) is pure omission --- fill four slots, leave the rest null:
\begin{lstlisting}
{"archetype": "Precision Optimization", "base_parent": "A",
 "justification": "the rigid 8-step structure is a bottleneck; simplify by removing steps",
 "slot_1": {"id": "a1", "kind": "keep"}, "slot_2": {"id": "a2", "kind": "keep"},
 "slot_3": {"id": "a3", "kind": "keep"}, "slot_4": {"id": "a4", "kind": "keep"},
 "slot_5": null, "slot_6": null, "slot_7": null, "slot_8": null}
\end{lstlisting}
The intervening hops (gen3--gen5) are all-\texttt{keep} diffs that re-run the same genome; the
fitness wobble across them (0.487, 0.421, 0.506) is executor sampling noise, not diff-driven
improvement. The honest reading: on this task and this weak mutator, most of the diff arm's
motion is structural bookkeeping riding a stochastic executor, with two genuine
grammar-legal structural edits (extend, prune) at the ends. What the diff \emph{guarantees}
is that every one of these six programs was well-formed by construction --- none of them could
have been the malformed JSON that \armA{} persisted 40 times.

\section{Verdict}
The hypothesis holds \emph{most} strongly exactly where it should: with a weak, cheap,
non-reasoning mutator, free-form JSON rewriting persisted \textbf{43 invalid children out of
100} (40 malformed JSON) and needed 39 pre-persist retries, while the constrained diff
persisted \textbf{0 invalid children out of 100} with a single retried pre-persist
schema rejection. The non-reasoning model also unmasks the diff's output economy that the strong-model
run had hidden: \textbf{diff output is 35\% of the rewrite's} (420 vs 1{,}204 tokens per
completed call) at near-parity input. Fitness is at parity at the top ($\sim$0.50), with a
higher diff child-mean (0.413 vs 0.362); the two operators diverge structurally --- the
rewrite collapses chains (its best is a 1-step degenerate), the diff preserves and locally
edits them (best a coherent 4-step prune at 0.525, above the rewrite's peak), and every diff
child's history reads as a clean series of typed edits (\S\ref{sec:series}). The one cost that
grows is downstream executor tokens, because the diff keeps longer chains alive. Caveats
unchanged: single replicate per arm, no same-config variance floor, \texttt{memory=none} so
citation-integrity is instrumented but trivially \texttt{0/0}, and \texttt{num\_parents=1} so
crossover was out of scope. Next: a variance floor on this task, a two-parent run to exercise
crossover diffs, and a memory-on run to give the evidence surface something to cite.
\end{document}
